\subsection{O Wisdom: the Great Antiphons}\label{sec:o-antiphons}

On the evening of 17 December the Church begins to call Christ by his names.

At Vespers she sings the Magnificat, the song Mary sang when she visited
Elizabeth: ``My soul doth magnify the Lord. And my spirit hath rejoiced in God
my Saviour.'' It ends with a promise kept: ``He hath received Israel his
servant, being mindful of his mercy. As he spoke to our fathers: to Abraham and
to his seed for ever.'' Before the canticle and after it the Church sings an
antiphon, said whole both times, and for seven evenings, to the day before the
Vigil, the antiphon is one of seven that begin with the same cry. Each calls
Christ by a name from the Scriptures, each recalls what he did or promised
under that name, and each ends with a single plea: \latin{veni}, come. So
Mary's song of the promise kept is sung each evening inside the Church's cry
for its coming.

\begin{description}

\item[17 December, \latin{O Sapientia}.] ``O Wisdom, That comest out of the
mouth of the Most High''. Wisdom says in Ecclesiasticus, ``I came out of the
mouth of the most High''. The first antiphon goes back before creation, to the
Word by whom all things were ordered, and asks him to come ``to teach us the
way of prudence''.

\item[18 December, \latin{O Adonai}.] ``O Adonai, and Ruler of the house of
Israël, Who didst appear unto Moses in the burning bush''. \latin{Adonai} is
Hebrew for ``Lord'', the word the Jews said in place of the Name of God, and
Exodus tells how ``the Lord appeared to him in a flame of fire out of the midst
of a bush''. The Child of Bethlehem is the God of the burning bush.

\item[19 December, \latin{O radix Jesse}.] ``O Root of Jesse, Which standest
for an ensign of the people''. Isaiah foretold ``the root of Jesse, who
standeth for an ensign of the people, him the Gentiles shall beseech''. From
the root of Jesse God raises a sign for all the nations, and the Church asks
him to ``make no tarrying!''

\item[20 December, \latin{O clavis David}.] ``O Key of David, and Sceptre of the
house of Israël; That openest, and no man shutteth''. Isaiah: ``I will lay the
key of the house of David upon his shoulder: and he shall open, and none shall
shut''. The antiphon asks him ``to bring out the prisoners from the prison''. The
prison is the shadow of death, and the key is in his hand.

\item[21 December, \latin{O Oriens}.] ``O Day-Spring, Brightness of the
everlasting Light, Sun of Righteousness''. Zachary sang of it in St Luke's
Gospel: ``the Orient from on high hath visited us''. \latin{Oriens} is the
rising of the sun, the east where the light first comes.

\item[22 December, \latin{O Rex Gentium}.] ``O King of the Gentiles, yea, and
Desire thereof, O Corner-stone that makest of twain one''. St Paul says of
Christ, ``he is our peace, who hath made both one''. The King of all the
nations is asked to save the creature he once shaped ``of the dust of the
earth''.

\item[23 December, \latin{O Emmanuel}.] ``O Emmanuel, our King and our
Law-giver''. Isaiah had said of the Virgin's Son, ``his name shall be called
Emmanuel'', and St Matthew tells what it means: ``God with us''. The
last antiphon calls him by that name, and asks him by it, for the last time, to
come.

\end{description}

The seven are a single prayer said over a week. They begin in eternity, with
Wisdom coming forth from the mouth of God, and they end at the edge of
Bethlehem, with God among us. In between they pass through the Old Testament:
Moses and the Law, the root of Jesse and the key of David, the dawn the
prophets promised and the King of the nations. Each night the Church calls him
by one more of his names, and each night she says the same word to him.

There is one more thing in them, and it lies in their letters. Read the Latin
names backward, from the last evening to the first, \latin{Emmanuel, Rex,
Oriens, Clavis, Radix, Adonai, Sapientia}, and their first letters spell
\latin{ero cras}: ``I shall be, tomorrow.'' No one can say whether anyone meant
it. The words are there all the same, and on the evening of 23 December, when
the last antiphon has been sung, the next day is the Vigil.

The seven are printed whole, in Latin and English, in
Appendix~\ref{sec:prayers}. Say the evening's antiphon at home, before the
Magnificat and after it if you pray Vespers, or alone if you do not; learn the
Latin of one of them. If a day is hard, take its antiphon into the next
morning, and let it be the first thing you say.
